\documentclass{article}
\usepackage[utf8]{inputenc}
\usepackage[T1]{fontenc}
\usepackage{geometry}
\geometry{
    a4paper,
    left=2.5cm,
    right=2.5cm,
    top=2.5cm,
    bottom=2.5cm
}
\usepackage{graphicx}
\usepackage{float}
\usepackage{hyperref}
\usepackage[english]{babel}
\usepackage{amsmath}
\usepackage{amssymb}
\usepackage{array}
\usepackage{booktabs}
\usepackage{caption}
\usepackage{subcaption}
\usepackage{enumitem}
\usepackage{fancyhdr}
\usepackage{titlesec}
\usepackage{url}
\usepackage{textcomp}
\usepackage{xcolor}
\usepackage{listings}
\usepackage{multirow}
\usepackage{longtable}
\usepackage{tabularx}
\usepackage{times}

% Set up page style
\pagestyle{fancy}
\fancyhf{}
\rfoot{\thepage}
\renewcommand{\headrulewidth}{0pt}

% Title formatting
\titleformat{\chapter}[display]{\normalfont\bfseries}{\chaptertitle\thechapter}{0pt}{\Huge}
\titlespacing*{\chapter}{0pt}{50pt}{40pt}

% Document info
\title{Template Document}
\author{Academic Support Agent}
\date{\today}

\begin{document}
\maketitle

\selectlanguage{brazil}
% \selectlanguage{english}

% Define o t tulo da tese.
\def\thesisTitle{T tulo da Disserta o de Mestrado} 

% Defina o nome do(a) autor(a) (  o seu nome).
\def\thesisAuthor{Nome do(a) Autor(a)} 

% Defina o nome do(a) orientador(a) 
\def\thesisAdvisor{Nome do(a) Orientador(a)} 

% Defina o nome do(a) co-orientador(a) 
% Comente a linha abaixo se n o tiver coorientador 
\def\thesisCoAdvisor{Nome do(a) Coorientador(a)} 

\pagestyle{empty}

% Content from cover/titulo would be included here
 

% Content from cover/ficha would be included here
 

% Content from cover/assinaturas would be included here
 

% Includes page number at the bottom of the page.
\pagestyle{plain}

% Uses small roman numbers for page numbering.
\pagenumbering{roman} 				

% Resets the page counter to 1.
\setcounter{page}{1}

% Content from cover/abstract would be included here

% Content from cover/resumo would be included here
 
% Content from cover/prefacio would be included here
 
% Content from cover/agradecimentos would be included here
 
% Content from cover/dedicatoria would be included here
 % opcional 


%Prints the table of contents
\renewcommand{\contentsname}{Sum rio}

%\addcontentsline{toc}{chapter}{Sum rio}


% Describes the research line and how the thesis is related with it 
% Content from cover/linha would be included here


% Prints a list of prizes and bibliographic and technical production 
% Content from cover/producao would be included here



% Prints a list of tables

\addcontentsline{toc}{chapter}{Lista de Tabelas}



% Prints a list of figures

\addcontentsline{toc}{chapter}{Lista de Figuras}


% Prints a list of abbreviations
% Content from cover/abreviacoes would be included here


% Prints a list of symbols
% Content from cover/simbolos would be included here



% Fancing heading alternating section and chapter with page count on the top
\pagestyle{fancy}
\fancyhead[LE,RO]{\thepage} 
\fancyhead[RE,LO]{\nouppercase{\leftmark}} 
\setlength{\headheight}{14.5pt}
\cfoot{}

\pagenumbering{arabic} 				
% Uses arabic numbers (normal numbers) for page numbering.

\setcounter{page}{1}
% Resets the page counter to 1.

% The chapters
% Content from chapters/a_introducao would be included here

% Content from chapters/b_revisao would be included here

% Content from chapters/c_metodologia would be included here

% Content from chapters/d_resultados would be included here

% Content from chapters/e_discussao would be included here

% Content from chapters/f_conclusoes would be included here


%Without the word "section"
\fancyhead[LO]{\nouppercase{\rightmark}} 

\renewcommand{\bibname}{Refer ncias} 
\addcontentsline{toc}{chapter}{Refer ncias}
 % note o uso do pacote natbib (para uso de referencia direta e indireta com \citet, \citep, etc.). 
%  % Caso a proposta seja em ingl s, descomentar para usar o estilo apalike original que vem com o pacote natbib. 


% To start appendices with different chapter numbering


\fancyhead[LO]{\nouppercase{\rightmark}} 

% Content from chapters/z_apendices.tex would be included here
\end{document}